\begin{corollary}[6.11.7]
\label{IV.2.s6.11.7-fr}
Con las notaciones de \textup{(6.11.6)}:
\begin{enumerate}
\item[\textup{(i)}] El conjunto $U_{S_1}(\mathcal F)$ es abierto en $X$.
\item[\textup{(ii)}] Para que el conjunto $U_{S_2}(\mathcal F)$ sea abierto, basta
que todo punto maximal $x$ de $\operatorname{Supp}(\mathcal F)$,
perteneciente a $U_{S_2}(\mathcal F)$, sea un punto interior de
$U_{S_2}(\mathcal F)$.
\end{enumerate}
\end{corollary}

\begin{proof}
La afirmación (i) se probó durante la demostración de (6.11.6);
por otra parte, para $n=2$, el caso $2^\circ$ de la demostración de (6.11.6)
es válido sin ninguna hipótesis sobre $X$, puesto que (con las mismas
notaciones) $U_{S_1}(\mathcal F)$ y
$U_{S_1}(\mathcal F/f\mathcal F)$ son abiertos en $X$. En cuanto al caso
$1^\circ$ de esa demostración, la hipótesis asegura precisamente que no es
necesario considerarlo.
\end{proof}
